\noindent\textbf{INTRODUÇÃO}
$\!$\\

Um tópico fundamental na gestão de investimentos é a otimização de portfólios financeiros, que busca equilibrar retorno e risco por meio da distribuição de recursos entre diferentes ativos. Com esse propósito, surgiram, na década de 1950, modelos como o de Média-Variância de Harry Markowitz, que introduziu o conceito de diversificação eficiente na teoria de portfólios. O modelo considera o retorno esperado dos ativos e a variância e covariância entre seus retornos como elementos fundamentais para a composição da carteira. Markowitz demonstrou que o desempenho de uma carteira não depende exclusivamente do retorno e do risco individual dos ativos, mas também da relação entre os ativos que a compõem, especialmente por meio de suas correlações \cite{markowitz}.

O modelo Black-Litterman (BL), introduzido por Fischer Black e Robert Litterman em 1992, oferece uma abordagem para a composição de portfólios que combina as expectativas implícitas no equilíbrio de mercado com visões específicas sobre os ativos. O modelo parte dos retornos de equilíbrio derivados das participações de mercado e incorpora visões dos investidores considerando diferentes níveis de confiança nessas estimativas. Fundamentado em uma formulação bayesiana, o modelo Black-Litterman permite combinar essas informações de maneira a reduzir a sensibilidade da carteira a estimativas isoladas de retorno \cite{blacklitterman}. Dessa forma, o modelo fornece uma estrutura para incorporar novas informações às expectativas de mercado sem desconsiderar o equilíbrio implícito nos preços dos ativos.

Além das informações provenientes do mercado, informações textuais, como notícias financeiras, podem fornecer sinais adicionais sobre as expectativas dos agentes em relação aos ativos. Nesse contexto, técnicas de Processamento de Linguagem Natural (PLN) e análise de sentimento podem ser utilizadas para identificar e quantificar informações subjetivas presentes em textos financeiros, permitindo sua incorporação em modelos quantitativos de análise e previsão \cite{du2024financial}. Paralelamente, técnicas de aprendizado de máquina (do inglês ML, \textit{Machine Learning}) podem ser empregadas na identificação de padrões nos dados históricos dos ativos, permitindo a construção de estimativas de retorno. Estudos anteriores exploraram, por exemplo, o emprego do XGBoost em modelos de previsão do mercado acionário \cite{shi2022attclx}.

Neste projeto, é proposta uma abordagem híbrida para a otimização de portfólios de ações brasileiras, combinando previsões obtidas por ML, informações de sentimento extraídas de notícias e o equilíbrio de mercado utilizado pelo modelo BL. O pipeline desenvolvido integra essas diferentes fontes de informação em etapas sequenciais: inicialmente, dados históricos dos ativos são utilizados para gerar previsões de retorno por meio do XGBoost; paralelamente, notícias financeiras são coletadas e submetidas à análise de sentimento; posteriormente, as informações quantitativas e qualitativas são combinadas na construção das visões utilizadas pelo modelo BL, considerando a confiança associada a cada fonte de informação. A partir dessas estimativas, são determinadas as alocações da carteira, que são posteriormente avaliadas por meio de \textit{backtesting} com dados históricos, considerando métricas de desempenho e risco e comparando os resultados com estratégias convencionais de otimização de portfólio. Em um experimento complementar, são comparados o XGBoost e outros modelos de previsão de retorno, utilizando os mesmos ativos, horizontes e períodos de avaliação.

Além da introdução, o trabalho está organizado em cinco capítulos. O primeiro capítulo apresenta a fundamentação teórica e aborda os seguintes temas: teoria moderna do portfólio, modelo Black-Litterman, aprendizado supervisionado e análise de sentimento em textos financeiros. O segundo capítulo descreve a arquitetura do sistema proposto, detalhando os componentes de previsão quantitativa, análise de sentimento e integração das informações para a otimização das carteiras. O terceiro capítulo apresenta o procedimento experimental, incluindo as bases de dados, o pré-processamento, as configurações dos modelos, a comparação complementar entre modelos de previsão de retorno, o protocolo de simulação e as métricas de avaliação. O quarto capítulo analisa e discute os resultados obtidos, considerando o desempenho das previsões e das estratégias de alocação, bem como as limitações do experimento. Por fim, o quinto e último capítulo reúne as conclusões, as principais contribuições do estudo e as sugestões para trabalhos futuros.

